% !TeX root = booklet-02.tex
\subsection{جریان‌های نشتی}
\label{sec:leakage}
در افزاره واقعی، بایاس خاموش نیز با جریان‌های کوچکی همراه است. کتاب
سه مسیر اصلی را معرفی می‌کند. \lr{Subthreshold Leakage} در مسیر سورس
و درین، \lr{Gate Leakage} از راه عایق گیت و \lr{Junction Leakage} در
پیوندهای ناحیه‌های \lr{Diffusion} و بدنه جریان دارند
\cite[بخش \secnum{۲.۴.۴} و شکل \secnum{۲.۱۹}]{weste}
\cite[اسلایدهای ۲۱ و ۲۲]{nonideal}. بایاس، دما، هندسه و ساختار فرایند،
سهم هر مسیر را تعیین می‌کنند.

در روایت فناوری کتاب سال ۲۰۱۱، نشتی در نمونه‌های با ابعاد بالاتر از
\(\qty{180}{\nano\metre}\) معمولاً سهم کوچکی داشته است. برای نمونه‌های
\(\qty{90}{\nano\metre}\) و \(\qty{65}{\nano\metre}\)، جریان‌های
زیرآستانه در مرتبه یک تا چند ده نانوآمپر برای هر ترانزیستور گزارش
شده‌اند. در نمونه‌های \(\qty{45}{\nano\metre}\)، نازک‌شدن عایق، نشتی
گیت را نیز مهم می‌کند و مرجع نقش عایق‌های \lr{High-k} را مطرح می‌کند
\cite[صفحه ۸۱]{weste}. این اعداد به نمونه‌ها و زمان مرجع مربوط‌اند.

\subsubsection{\lr{Subthreshold Leakage}}
\label{sec:subthreshold-leakage}
\paragraph*{وارونگی ضعیف و معنای خاموشی}
در \lr{Weak Inversion}، ولتاژ گیت زیر آستانه است و مقدار جریان به‌صورت
نمایی با ولتاژ گیت تغییر می‌کند. در توصیف کتاب، عبور حرارتی حامل‌ها
از سد پتانسیل، جریان زیرآستانه را ایجاد می‌کند
\cite[بخش \secnum{۲.۴.۴.۱}]{weste}. در مدل سوئیچی، خروجی منطقی بر
اساس مسیرهای روشن و خاموش تحلیل می‌شود. مدل نشتی، مقدار جریان کوچک
مسیر خاموش و اثر آن بر توان و نگهداری بار را نیز محاسبه می‌کند.

\paragraph*{پیش‌نیاز تابع نمایی و لگاریتم ده‌دهی}
عدد \(e\approx2.71828\) پایه لگاریتم طبیعی است. تابع
\(\exp(x)=e^x\) برای هر \(x\) مثبت است. برای نمونه،
\(e^0=1\)، \(e^1\approx2.718\)، \(e^{-1}=1/e\approx0.368\) و
\(e^{-3}\approx0.0498\). منفی‌ترشدن توان، مقدار تابع را کوچک‌تر می‌کند.
قانون \(e^{a+b}=e^ae^b\) و در نتیجه \(e^a/e^b=e^{a-b}\) در محاسبه
نسبت نشتی‌ها به کار می‌روند.

لگاریتم ده‌دهی \(\log_{10}x\)، توانی از ده است که عدد \(x>0\) را
می‌سازد. مثلاً \(\log_{10}1000=3\) و \(\log_{10}0.01=-2\).
یک \lr{Decade} تغییر جریان یعنی ضرب یا تقسیم آن در ده.
برای تبدیل پایه از رابطه \(\log_{10}x=\ln x/\ln10\) استفاده می‌کنیم.
همچنین \(e^x=10^{x/\ln10}\) است. نسبت‌های داخل لگاریتم و توان تابع
نمایی باید بی‌بعد باشند.

\paragraph*{مدل جریان و تعریف همه پارامترها}
رابطه \secnum{۲.۴۲} کتاب و اسلاید ۲۳ برای \lr{nMOS} با
\(V_{DS}\ge0\) چنین بازنویسی می‌شود
\begin{equation}
 I_{\mathrm{sub}}=I_{ds0}
 \exp\!\left(\frac{V_{GS}-V_{t0}+\eta V_{DS}-k_\gamma V_{SB}}{nv_T}\right)
 \left(1-e^{-V_{DS}/v_T}\right).
 \label{eq:leakage-subthreshold-model}
\end{equation}
در این رابطه، \(V_{t0}\) آستانه مرجع، \(\eta\) ضریب \lr{DIBL} و
\(k_\gamma\) ضریب خطی اثر بدنه از
\eqref{eq:body-linearization} هستند. \(n>0\) ضریب وابسته به فرایند
است که مرجع مقدار معمول آن را حدود \num{1.3} تا \num{1.7} می‌گیرد.
\(v_T=kT/q\) ولتاژ حرارتی است. \(I_{ds0}\) ضریب جریان وابسته به
فرایند و هندسه است \cite[صفحات ۸۱ و ۸۲]{weste}.

صورت کسر داخل نمایی، یک ولتاژ است. مخرج \(nv_T\) نیز ولتاژ است، پس
توان بی‌بعد می‌شود. عامل آخر هم بی‌بعد است و واحد کل جریان از
\(I_{ds0}\) به دست می‌آید. برای استفاده از آستانه مؤثر، می‌توان
\(V_{t,\mathrm{eff}}=V_{t0}-\eta V_{DS}+k_\gamma V_{SB}\) را تعریف کرد
و صورت توان را \(V_{GS}-V_{t,\mathrm{eff}}\) نوشت.
این صورت از تقریب خطی بایاس بدنه استفاده می‌کند و دامنه آن با
\eqref{eq:body-linear-error} سنجیده می‌شود.

مدل مرجع، ضریب جریان را با رابطه تجربی زیر نیز تقریب می‌زند
\begin{equation}
 I_{ds0}=\beta v_T^2e^{1.8},\qquad
 \beta=\mu C_{\mathrm{ox}}\frac{W}{L}.
 \label{eq:leakage-current-prefactor}
\end{equation}
عدد \(1.8\) ضریب تجربی همان مدل است
\cite[رابطه \secnum{۲.۴۳}]{weste}. حاصل‌ضرب
\(\beta\) با واحد \(\unit{\ampere\per\volt\squared}\) و \(v_T^2\)،
واحد آمپر دارد. در شرایطی که صورت توان صفر و
\(V_{DS}\gg v_T\) باشد، جریان به \(I_{ds0}\) نزدیک می‌شود.
پارامترهای این رابطه می‌توانند با دما و فرایند تغییر کنند.

\paragraph*{عامل درین و معیار تقریب آن}
برای بررسی عامل آخر، \(D=V_{DS}\) و
\(F(D)=1-e^{-D/v_T}\) را تعریف می‌کنیم. در \(D=0\)، عامل \(F\)
صفر و جریان مدل نیز صفر است. با افزایش \(D>0\)، این عامل به یک
نزدیک می‌شود. چند مقدار در جدول \ref{tab:leakage-drain-factor}
آمده‌اند.
\begin{table}[!htbp]
\centering
\small
\caption{عامل درین در مدل زیرآستانه و فاصله نسبی آن از یک.}
\label{tab:leakage-drain-factor}
\begin{tabular}{cccc}
\toprule
\(D/v_T\) & \(e^{-D/v_T}\) & \(F(D)\) & فاصله از یک، درصد\\
\midrule
\num{1} & \num{0.3679} & \num{0.6321} & \num{36.79}\\
\num{2} & \num{0.1353} & \num{0.8647} & \num{13.53}\\
\num{3} & \num{0.0498} & \num{0.9502} & \num{4.98}\\
\num{5} & \num{0.00674} & \num{0.99326} & \num{0.674}\\
\bottomrule
\end{tabular}
\end{table}

برای فاصله حداکثر ۱ درصد از یک، باید
\(e^{-D/v_T}\le0.01\) باشد. با لگاریتم‌گیری،
\(D\ge-v_T\ln0.01\approx4.605v_T\) به دست می‌آید.
در \(v_T=\qty{26}{\milli\volt}\)، این شرط حدود
\(D\ge\qty{120}{\milli\volt}\) است. در
\(D=\qty{50}{\milli\volt}\)، عامل برابر \(0.8538\) و در
\(D=\qty{100}{\milli\volt}\)، برابر \(0.9786\) است.
در استفاده از تعبیر «چند برابر ولتاژ حرارتی» در متن مرجع، دقت مورد
نیاز محاسبه تعیین‌کننده انتخاب تقریب است.

در \(D\ll v_T\)، تقریب \(e^{-D/v_T}\approx1-D/v_T\) می‌دهد
\(F(D)\approx D/v_T\). ازاین‌رو، جریان در نزدیکی درین صفر با
\(D\) متناسب می‌شود. در \(D\gg v_T\)، عامل \(F\) تقریباً ثابت
است، اما جمله \(\eta D\) در توان هنوز می‌تواند با افزایش درین،
نشتی را افزایش دهد.

\paragraph*{مثال آموزشی مقدار جریان و اثر گیت}
فرض کنید \(\beta=\qty{1}{\milli\ampere\per\volt\squared}\)،
\(n=1.4\)، \(v_T=\qty{0.026}{\volt}\)،
\(V_{t0}=\qty{0.35}{\volt}\)، \(\eta=0.05\)،
\(V_{DS}=\qty{1}{\volt}\) و \(V_{SB}=0\) باشند.
این پارامترها برای یک مثال آموزشی فرض شده‌اند. ابتدا داریم
\begin{equation}
 I_{ds0}=(10^{-3})(0.026)^2e^{1.8}\,\unit{\ampere}
 \approx\qty{4.090}{\micro\ampere}.
 \label{eq:leakage-prefactor-example}
\end{equation}
در \(V_{GS}=0\)، صورت توان
\(-0.35+0.05(1)=-\qty{0.30}{\volt}\) و مخرج
\(1.4(0.026)=\qty{0.0364}{\volt}\) است.
چون \(F(\qty{1}{\volt})\approx1\)، جریان برابر است با
\begin{equation}
 I_{\mathrm{off}}\approx(4.090\times10^{-6})e^{-0.30/0.0364}
 \approx\qty{1.077}{\nano\ampere}.
 \label{eq:leakage-absolute-example}
\end{equation}
اگر فقط گیت به \(\qty{0.1}{\volt}\) برسد، نسبت دو جریان از قانون
تقسیم توابع نمایی به دست می‌آید
\begin{equation}
 \frac{I(\qty{0.1}{\volt})}{I(0)}
 =e^{0.1/0.0364}\approx15.60,
 \qquad I(\qty{0.1}{\volt})\approx\qty{16.81}{\nano\ampere}.
 \label{eq:leakage-gate-ratio-example}
\end{equation}
آستانه مؤثر در هر دو نقطه \(\qty{0.30}{\volt}\) است و هر دو بایاس
گیت در دامنه زیرآستانه قرار دارند.

\paragraph*{استخراج \lr{Subthreshold Slope} و خواندن نمودار لگاریتمی}
در بایاس درین، بدنه و دمای ثابت، تمام عامل‌های مستقل از گیت را در یک
ضریب ثابت جمع می‌کنیم. برای یک جریان مرجع ثابت \(I_r>0\)، از
\eqref{eq:leakage-subthreshold-model} نتیجه می‌شود
\begin{equation}
 y=\log_{10}\!\left(\frac{I_{\mathrm{sub}}}{I_r}\right)
 =c+\frac{V_{GS}}{nv_T\ln10}.
 \label{eq:leakage-log-line}
\end{equation}
نماد \(c\) یک مقدار ثابت در این مقایسه است. پس نمودار \(y\) برحسب
\(V_{GS}\) خطی است و شیب آن \(1/(nv_T\ln10)\) می‌شود.
کتاب معکوس این شیب را \lr{Subthreshold Slope} می‌نامد
\begin{equation}
 \boxed{S=\left[\frac{\mathrm{d}\log_{10}(I_{\mathrm{sub}}/I_r)}
 {\mathrm{d}V_{GS}}\right]^{-1}=nv_T\ln10}.
 \label{eq:leakage-subthreshold-slope}
\end{equation}
این رابطه با \secnum{۲.۴۴} کتاب و اسلاید ۲۳ متناظر است
\cite{weste,nonideal}. \(S\) تغییر گیت لازم برای ده‌برابرشدن جریان
است و معمولاً با واحد \(\unit{\milli\volt}\) به ازای یک
\lr{Decade} بیان می‌شود. در مثال قبل،
\(S=1.4(0.026)\ln10\approx\qty{83.81}{\milli\volt}\) است.
افزایش گیت به اندازه \(2S\)، جریان را صد برابر می‌کند.

شکل \ref{fig:leakage-semilog} نسبت جریان به مقدار گیت صفر را برای
همین \(S\) نشان می‌دهد. فاصله برابر روی محور عمودی لگاریتمی، به معنای
ضریب برابر در جریان است. برای \(n\) بین \num{1.3} و \num{1.7} و
\(v_T=\qty{26}{\milli\volt}\)، مقدار \(S\) حدود
\(\qty{77.8}{\milli\volt}\) تا \(\qty{101.8}{\milli\volt}\) به ازای هر دهه است.
\begin{figure}[!htbp]
\centering
\begin{latin}
\PersianMathDigits
\begin{tikzpicture}[x=22cm,y=1.05cm]
 \draw[->] (-0.12,-1.5)--(0.24,-1.5) node[right] {\(V_{GS}\;(\mathrm{V})\)};
 \draw[->] (-0.12,-1.5)--(-0.12,2.8) node[above] {\(I/I_{\mathrm{off}}\)};
 \foreach \x in {-0.1,0,0.1,0.2} {\draw (\x,-1.5)--(\x,-1.56) node[below] {\(\x\)};}
 \foreach \y/\value in {-1/0.1,0/1,1/10,2/100} {
  \draw (-0.12,\y)--(-0.125,\y) node[left] {\(\value\)};
 }
 \draw[thick,domain=-0.1:0.2,samples=50] plot (\x,{\x/0.0838141});
 \draw[dotted] (0,0)--(0.0838141,0)--(0.0838141,1)--(0,1);
 \draw[<->] (0,0.15)--(0.0838141,0.15) node[midway,above] {\(S\)};
\end{tikzpicture}
\end{latin}
\caption{نمودار آموزشی زیرآستانه با محور عمودی لگاریتمی و
\(S=\qty{83.81}{\milli\volt}\) به ازای هر دهه. گیت به اندازه \(S\)
تغییر می‌کند و نسبت جریان از یک به ده می‌رسد. بایاس‌های دیگر ثابت‌اند.}
\label{fig:leakage-semilog}
\end{figure}

در شکل \secnum{۲.۲۰} کتاب، محور عمودی جریان بر واحد عرض با واحد
\(\unit{\ampere\per\micro\metre}\) است. برای نمونه
\(\qty{65}{\nano\metre}\) در دمای \(\qty{70}{\degreeCelsius}\)،
مقادیر حدود \(\qty{27}{\nano\ampere\per\micro\metre}\) و
\(\qty{2}{\nano\ampere\per\micro\metre}\) در گیت صفر و دو بایاس
درین \(\qty{1.0}{\volt}\) و \(\qty{0.1}{\volt}\) مشخص شده‌اند.
نسبت آن‌ها حدود \(13.5\) است. بخش تقریباً خطی منحنی زیرآستانه و
بالارفتن جریان در گیت بسیار منفی، به‌ترتیب رفتار نمایی و اثر
\lr{GIDL} را نشان می‌دهند \cite[صفحه ۸۱]{weste}.

\paragraph*{بازنویسی دقیق نسبت به \(I_{\mathrm{off}}\)}
جریان مرجع را در گیت و بایاس بدنه صفر و درین \(V_{DD}>0\) تعریف
می‌کنیم. از \eqref{eq:leakage-subthreshold-model} داریم
\begin{equation}
 I_{\mathrm{off}}=I_{ds0}
 e^{(-V_{t0}+\eta V_{DD})/(nv_T)}F(V_{DD}).
 \label{eq:leakage-off-definition}
\end{equation}
تقسیم رابطه جریان بر رابطه مرجع، عامل \(I_{ds0}\) و آستانه مرجع را
حذف می‌کند. با استفاده از \(S=nv_T\ln10\)، نتیجه برابر است با
\begin{equation}
 \boxed{I_{\mathrm{sub}}=I_{\mathrm{off}}
 10^{[V_{GS}+\eta(V_{DS}-V_{DD})-k_\gamma V_{SB}]/S}
 \frac{F(V_{DS})}{F(V_{DD})}}.
 \label{eq:leakage-off-normalized}
\end{equation}
در بایاس مرجع، توان صفر و نسبت دو عامل \(F\) برابر یک است و جریان
مرجع دقیقاً بازیابی می‌شود. در بازنویسی رابطه \secnum{۲.۴۵} کتاب،
عامل مخرج با فرض \(V_{DD}\gg v_T\) تقریباً یک گرفته شده است.
برای نگه‌داشتن همه عامل‌ها در بایاس درین کوچک، از
\eqref{eq:leakage-off-normalized} استفاده می‌کنیم.

این نسبت برای هندسه، دما و پارامترهای یکسان نوشته شده است.
در یک دمای ثابت، افزایش \(V_{t0}\) به اندازه \(\Delta V_t>0\)،
جریان را در عامل \(e^{-\Delta V_t/(nv_T)}\) ضرب می‌کند.
افزایش بایاس بدنه در تقریب خطی، عامل
\(e^{-k_\gamma\Delta V_{SB}/(nv_T)}\) و تغییر درین، هم عامل نمایی
\lr{DIBL} و هم نسبت دو \(F\) را وارد می‌کند.
برای \lr{pMOS} با قرارداد بزرگی‌های مثبت، همین ساختار با
\(V_{SG},V_{SD},V_{BS}\)، بزرگی آستانه و پارامترهای آن افزاره به کار
می‌رود.

\paragraph*{حل مثال \secnum{۲.۶} کتاب و نسبت هزاربرابری}
مثال مرجع می‌خواهد جریان گیت صفر، هزار برابر کوچک‌تر از جریان در
گیت برابر آستانه باشد. در یک بایاس درین ثابت با \(V_{DS}\gg v_T\)،
بایاس بدنه صفر و با قرارداد آستانه همان مقایسه، نسبت جریان‌ها چنین است
\begin{equation}
 \frac{I(0)}{I(V_t)}=e^{-V_t/(nv_T)}=10^{-3}.
 \label{eq:leakage-three-decades}
\end{equation}
از دو طرف لگاریتم می‌گیریم. با \(\ln10^{-3}=-3\ln10\)، داریم
\begin{equation}
 -\frac{V_t}{nv_T}=-3\ln10
 \quad\Longrightarrow\quad V_t=3nv_T\ln10=3S.
 \label{eq:leakage-minimum-threshold}
\end{equation}
پس برای \(n=1.5\) و \(v_T=\qty{26}{\milli\volt}\)،
\(S\approx\qty{89.8}{\milli\volt}\) و
\(V_t\approx\qty{269.4}{\milli\volt}\) است. گردکردن مرجع حدود
\(\qty{270}{\milli\volt}\) می‌دهد. برای نمونه \lr{Silicon-on-Insulator
(SOI)} مثال با \(n=1.3\)، مقدار \(S\approx\qty{77.8}{\milli\volt}\)
و \(V_t\approx\qty{233.5}{\milli\volt}\) می‌شود
\cite[مثال \secnum{۲.۶}، صفحه ۸۲]{weste}.
شرط حداقل آستانه برای نسبت حداکثر \(10^{-3}\)، همان
\(V_t\ge3S\) است.

\paragraph*{اثر بر توان و نگهداری بار}
توان ایستای دریافت‌شده از یک منبع ثابت از
\(P_{\mathrm{static}}=V_{DD}I_{\mathrm{supply}}\) به دست می‌آید.
در یک مثال فرضی با \(10^8\) ترانزیستور، کسر خاموش \(0.8\)، نشتی
متوسط هر ترانزیستور خاموش \(\qty{1}{\nano\ampere}\) و تغذیه
\(\qty{1}{\volt}\)، اگر این جریان‌ها سهم جمع‌پذیر همان منبع باشند،
جریان کل \(\qty{0.08}{\ampere}\) و توان \(\qty{80}{\milli\watt}\)
می‌شود. برای مدار مشخص، اتصال مسیرها و جریان دریافتی از هر منبع با
\lr{KCL} محاسبه می‌شود.

نشتی همچنین بار گره خازنی را در طول زمان تغییر می‌دهد. با جریان
تقریباً ثابت \(I\)، تغییر ولتاژ \(\Delta V\) روی ظرفیت \(C\) در زمان
\(t\approx C\Delta V/I\) رخ می‌دهد. مثلاً
\(C=\qty{10}{\femto\farad}\)، \(I=\qty{1}{\nano\ampere}\) و
\(\Delta V=\qty{0.2}{\volt}\)، زمان حدود \(\qty{2}{\micro\second}\)
می‌دهند. این برآورد آموزشی، حساسیت گره‌های ذخیره بار و نیاز به
\lr{Refresh} در حافظه پویا را توضیح می‌دهد
\cite[صفحه ۸۲]{weste}.

\subsubsection{\lr{Gate Leakage}}
\label{sec:gate-leakage}
\paragraph*{تونل‌زنی و دو سازوکار مرجع}
در عایق بسیار نازک، احتمال عبور کوانتومی حامل از سد عایق، جریان
تونل‌زنی ایجاد می‌کند. برای بیان مقیاس، کتاب عایق‌های نازک‌تر از حدود
\(\qty{1.5}{\nano\metre}\) تا \(\qty{2.0}{\nano\metre}\) را در این
بحث مطرح می‌کند. کتاب دو سازوکار \lr{Fowler-Nordheim (FN)
Tunneling} و \lr{Direct Tunneling} را معرفی می‌کند.
اولی در ولتاژ زیاد و ضخامت متوسط اهمیت پیدا می‌کند و در برنامه‌ریزی
\lr{EEPROM} به کار می‌رود. دومی در عایق نازک و ولتاژ پایین‌تر اهمیت
دارد \cite[بخش \secnum{۲.۴.۴.۲}]{weste}.

برای \lr{Direct Tunneling}، رابطه \secnum{۲.۴۷} کتاب و اسلاید ۲۴ را
با نام‌گذاری ضرایب مستقل می‌نویسیم
\begin{equation}
 I_{\mathrm{gate}}=WA_g
 \left(\frac{V_{DD}}{t_{\mathrm{ox}}}\right)^2
 \exp\!\left(-B_g\frac{t_{\mathrm{ox}}}{V_{DD}}\right).
 \label{eq:gate-tunneling-model}
\end{equation}
\(A_g,B_g\) همان ثابت‌های فناوری \(A,B\) در مرجع هستند.
زیرنویس‌ها برای تشخیص آن‌ها از ضرایب مدل پیوند به کار رفته‌اند.
\(W\) عرض و \(t_{\mathrm{ox}}\) ضخامت فیزیکی عایق در مدل است.
این مدل تجربی در شرایط بایاس و پارامترهای مربوط به برازش استفاده
می‌شود \cite[صفحه ۸۳]{weste}\cite[اسلاید ۲۴]{nonideal}.

با دستگاه طول متر، بی‌بعدبودن توان ایجاب می‌کند
\([B_g]=\unit{\volt\per\metre}\). چون
\(W(V_{DD}/t_{\mathrm{ox}})^2\) واحد
\(\mathrm{V}^2/\mathrm{m}\) دارد،
\([A_g]=\mathrm{A\,m\,V^{-2}}\) است. اگر ضرایب با نانومتر برازش
شده‌اند، ضخامت و واحد ضرایب نیز با همان قرارداد وارد می‌شوند.

\paragraph*{حساسیت ضخامت و مثال نسبت جریان}
در ولتاژ، عرض و ضرایب ثابت، برای دو ضخامت مثبت \(t_1,t_2\)، تقسیم
دو رابطه \eqref{eq:gate-tunneling-model} می‌دهد
\begin{equation}
 \frac{I_2}{I_1}=\left(\frac{t_1}{t_2}\right)^2
 \exp\!\left[\frac{B_g(t_1-t_2)}{V_{DD}}\right].
 \label{eq:gate-thickness-ratio}
\end{equation}
کاهش ضخامت، هم عامل مربعی و هم عامل نمایی را افزایش می‌دهد.
در مثال فرضی \(t_1=\qty{1.1}{\nano\metre}\)،
\(t_2=\qty{1.0}{\nano\metre}\)، \(V_{DD}=\qty{1}{\volt}\) و
\(B_g=\qty{10}{\volt\per\nano\metre}\)، نسبت برابر
\((1.1/1.0)^2e^1\approx3.29\) است. مقدار \(B_g\) در این مثال
آموزشی فرض شده است. برای جریان مطلق، مقدار \(A_g\) نیز لازم است.
با ثابت بودن دیگر پارامترها، دوبرابرشدن عرض در همین مدل جریان را
دو برابر می‌کند.

در متن مرجع، برای داده‌های مورد بحث، افزایش جریان حداقل با ضریب
حدود \(2.7\) به ازای کاهش یک آنگستروم ضخامت گزارش شده است.
کتاب مقیاس ضخامت \(\qty{1.05}{\nano\metre}\)، معادل \num{10.5} آنگستروم، و چگالی جریان
\(\qty{100}{\ampere\per\centi\metre\squared}\) و ضخامت حدود
\(\qty{0.3}{\nano\metre}\) برای هر لایه اتمی \lr{SiO2} را در توضیح
نمونه‌های همان دوره مطرح می‌کند \cite[صفحه ۸۳]{weste}.
این اعداد همراه ماده و شرایط نمودار مرجع تفسیر می‌شوند.

\paragraph*{ولتاژ، چگالی جریان و عایق \lr{High-k}}
در ضخامت ثابت، نسبت جریان در دو ولتاژ مثبت چنین است
\begin{equation}
 \frac{I(V_2)}{I(V_1)}=
 \left(\frac{V_2}{V_1}\right)^2
 \exp\!\left[B_gt_{\mathrm{ox}}
 \left(\frac{1}{V_1}-\frac{1}{V_2}\right)\right].
 \label{eq:gate-voltage-ratio}
\end{equation}
برای \(V_2>V_1\)، هر دو عامل نسبت را افزایش می‌دهند.
در شکل \secnum{۲.۲۱} کتاب، محور عمودی چگالی جریان گیت
\(J_G=I_{\mathrm{gate}}/A_{\mathrm{gate}}\) با واحد
\(\unit{\ampere\per\centi\metre\squared}\) است.
برای چگالی یکنواخت، جریان از ضرب چگالی در مساحت به دست می‌آید.
مثلاً \(J_G=\qty{1}{\ampere\per\centi\metre\squared}\) و مساحت
فرضی \(\qty{0.005}{\micro\metre\squared}
=\qty{5e-11}{\centi\metre\squared}\)، جریان
\(\qty{50}{\pico\ampere}\) می‌دهند.
در شکل زیرآستانه \secnum{۲.۲۰}، نرمال‌سازی با عرض انجام شده است.
واحد محور هر شکل، نحوه تبدیل به جریان کل را تعیین می‌کند.

برای جریان روشن زیاد، ظرفیت اکسید بر واحد سطح بزرگ مطلوب است.
از \(C_{\mathrm{ox}}=\varepsilon/t\)، افزایش \lr{Permittivity} عایق
امکان استفاده از ضخامت فیزیکی بیشتر را در ظرفیت یکسان فراهم می‌کند.
ضخامت معادل اکسید از \(\mathrm{EOT}=t\varepsilon_{\mathrm{ox}}/
\varepsilon\) به دست می‌آید. تونل‌زنی به ضخامت فیزیکی و ویژگی سد
ماده وابسته است، پس ضرایب نشتی نیز باید با ماده عایق سازگار باشند.
در نمونه‌های \lr{SiO2} مورد بحث کتاب، نشتی گیت \lr{nMOS} می‌تواند
در مرتبه ده برابر \lr{pMOS} باشد، زیرا حامل‌ها سدهای انرژی متفاوتی
دارند. مرجع وابستگی این رفتار به نوع عایق را یادآور می‌شود
\cite[صفحه ۸۳]{weste}.

\subsubsection{\lr{Junction Leakage}}
\label{sec:junction-leakage}
\paragraph*{مدل دیود و علامت ولتاژ پیوند}
پیوند ناحیه‌های \lr{Diffusion} با بدنه یا \lr{Well} و پیوند
\lr{Well} با \lr{Substrate} مانند دیود عمل می‌کنند.
در بایاس معمول، بدنه و \lr{Well} به زمین یا تغذیه متصل می‌شوند تا
پیوندها در بایاس صفر یا معکوس قرار گیرند
\cite[بخش \secnum{۲.۴.۴.۳} و شکل \secnum{۲.۲۲}]{weste}.
مدل دیود کتاب را با نماد \(I_{\mathrm{dio}}\) برای جریان و
\(I_{\mathrm{rev}}>0\) برای ضریب جریان معکوس بازنویسی می‌کنیم
\begin{equation}
 I_{\mathrm{dio}}=I_{\mathrm{rev}}
 \left(e^{V_{\mathrm{dio}}/v_T}-1\right),\qquad
 V_{\mathrm{dio}}=V_{\mathrm{anode}}-V_{\mathrm{cathode}}.
 \label{eq:junction-diode-model}
\end{equation}
این دو جریان با \(I_D,I_S\) رابطه \secnum{۲.۴۸} کتاب متناظرند.
نمادهای مرجع در این رابطه به جریان دیود و ضریب جریان معکوس اشاره
دارند \cite[صفحه ۸۴ و پانوشت آن]{weste}\cite[اسلاید ۲۶]{nonideal}.

در پیوند \lr{nMOS}، بدنه \lr{p-type} آند و ناحیه
\lr{n-type Diffusion} کاتد است. بنابراین برای پیوند درین و بدنه،
\(V_{\mathrm{dio}}=V_B-V_D=-V_{DB}\) است.
اگر \(V_{DB}>0\)، ولتاژ دیود منفی و جریان آن در قرارداد آند به
کاتد منفی می‌شود. در بایاس معکوس بزرگ‌تر از چند برابر \(v_T\)،
جمله نمایی به صفر نزدیک و
\(I_{\mathrm{dio}}\approx-I_{\mathrm{rev}}\) است.
جریان معکوس در این قرارداد از ناحیه درین به بدنه می‌رود.
در بایاس صفر نیز \(I_{\mathrm{dio}}=0\) حاصل می‌شود.

ضریب \(I_{\mathrm{rev}}\) به ناخالصی، مساحت و محیط ناحیه
\lr{Diffusion} وابسته است. کتاب برای مؤلفه معمول دیود، مرتبه
\(\qty{0.01}{\femto\ampere\per\micro\metre\squared}\) تا
\(\qty{0.1}{\femto\ampere\per\micro\metre\squared}\) را در
نمونه‌های مورد بحث ذکر می‌کند \cite[صفحه ۸۴]{weste}.
در مثال فرضی با چگالی \(\qty{0.1}{\femto\ampere\per\micro\metre\squared}\)،
مساحت \(\qty{0.2}{\micro\metre\squared}\) و صرف‌نظر از سهم محیط،
\(I_{\mathrm{rev}}=\qty{0.02}{\femto\ampere}\) می‌شود.
در بایاس معکوس \(\qty{1}{\volt}\) و
\(v_T=\qty{0.026}{\volt}\)، جریان حدود
\(-\qty{0.02}{\femto\ampere}\) است.

\paragraph*{\lr{Band-to-Band Tunneling} و نقش ناخالصی}
در پیوندهای با ناخالصی زیاد، میدان قوی می‌تواند
\lr{Band-to-Band Tunneling (BTBT)} ایجاد کند.
مرجع اثر \lr{Halo Doping}، یعنی افزایش ناخالصی نزدیک انتهای کانال
برای کنترل آستانه، را بر افزایش این نشتی توضیح می‌دهد.
وجود تله‌های ناشی از نقص ماده نیز می‌تواند
\lr{Trap-Assisted Tunneling (TAT)} را تقویت کند.
در توصیف کتاب، سهم مهمی از این جریان نزدیک دیواره پیوند رو به
کانال رخ می‌دهد \cite[صفحه ۸۴]{weste}\cite[اسلاید ۲۷]{nonideal}.

رابطه تجربی \secnum{۲.۴۹} کتاب برای بزرگی جریان چنین است
\begin{equation}
 I_{\mathrm{BTBT}}=WX_j A_b\frac{E_j}{E_g^{1/2}}V_{DD}
 \exp\!\left(-B_b\frac{E_g^{3/2}}{E_j}\right).
 \label{eq:junction-btbt-model}
\end{equation}
\(X_j\) عمق پیوند، \(E_j\) بزرگی میدان پیوند و \(A_b,B_b\)
ثابت‌های فناوری هستند. کتاب \(E_g\) را ولتاژ متناظر با گاف انرژی
تعریف می‌کند. اگر انرژی گاف برحسب ژول معلوم باشد، تقسیم آن بر
\(q\) مقدار \(E_g\) را برحسب ولت می‌دهد.
توان‌های \(1/2\) و \(3/2\) روی همین کمیت اعمال می‌شوند.
در قرارداد طول متر، بی‌بعدبودن توان و واحد آمپر نتیجه می‌دهند
\begin{equation}
 [B_b]=\mathrm{m}^{-1}\mathrm{V}^{-1/2},\qquad
 [A_b]=\mathrm{A\,m}^{-1}\mathrm{V}^{-3/2}.
 \label{eq:junction-btbt-units}
\end{equation}
ضرایب \(A_b,B_b\) با \(A_g,B_g\) مدل گیت جداگانه تعریف و برازش
می‌شوند.

\paragraph*{میدان پیوند و یک محاسبه آموزشی}
در بایاس معکوس \(V_{DD}\)، رابطه \secnum{۲.۵۰} کتاب میدان را با
ناخالصی \lr{Halo} و \lr{Source/Drain} مرتبط می‌کند
\begin{equation}
 E_j=\sqrt{\frac{2qN_{\mathrm{halo}}N_{\mathrm{sd}}}
 {\varepsilon_{\mathrm{si}}(N_{\mathrm{halo}}+N_{\mathrm{sd}})}
 \left[V_{DD}+v_T\ln\!\left(
 \frac{N_{\mathrm{halo}}N_{\mathrm{sd}}}{n_i^2}\right)\right]}.
 \label{eq:junction-field-model}
\end{equation}
همه غلظت‌ها باید با واحد یکسان و \(\varepsilon_{\mathrm{si}}\) با
واحد طول سازگار وارد شوند. برای دنبال‌کردن محاسبه،
دو کمیت کمکی را تعریف می‌کنیم
\begin{equation}
 N_{\mathrm{eff}}=\frac{N_{\mathrm{halo}}N_{\mathrm{sd}}}
 {N_{\mathrm{halo}}+N_{\mathrm{sd}}},\qquad
 \phi_{\mathrm{bi}}=v_T\ln\!\left(
 \frac{N_{\mathrm{halo}}N_{\mathrm{sd}}}{n_i^2}\right).
 \label{eq:junction-field-parameters}
\end{equation}
با جایگذاری آن‌ها، رابطه میدان به شکل ساده‌تر زیر تبدیل می‌شود
\[
 E_j=\sqrt{\frac{2qN_{\mathrm{eff}}}{\varepsilon_{\mathrm{si}}}
 (V_{DD}+\phi_{\mathrm{bi}})}.
\]
\(\phi_{\mathrm{bi}}\) واحد ولت دارد و عبارت زیر رادیکال، واحد
\(\mathrm{V}^2/\mathrm{m}^2\) یا شکل سازگار سانتی‌متری آن را دارد.

در مثال فرضی، \(N_{\mathrm{halo}}=\qty{e18}{\per\centi\metre\cubed}\)،
\(N_{\mathrm{sd}}=\qty{e20}{\per\centi\metre\cubed}\)،
\(n_i=\qty{1.45e10}{\per\centi\metre\cubed}\)،
\(v_T=\qty{0.026}{\volt}\) و
\(\varepsilon_{\mathrm{si}}=11.7(\qty{8.85e-14}{\farad\per\centi\metre})\)
هستند. محاسبه می‌دهد
\begin{align}
 N_{\mathrm{eff}}&\approx\qty{9.901e17}{\per\centi\metre\cubed},
 &\phi_{\mathrm{bi}}&\approx\qty{1.0583}{\volt},\notag\\
 E_j(V_{DD}=\qty{0.8}{\volt})&\approx\qty{7.541e5}{\volt\per\centi\metre},
 &E_j(V_{DD}=\qty{1.2}{\volt})&\approx\qty{8.313e5}{\volt\per\centi\metre}.
 \label{eq:junction-field-example}
\end{align}
برای مشاهده حساسیت نمایی، حاصل‌ضرب فرضی
\(B_bE_g^{3/2}=\qty{3e6}{\volt\per\centi\metre}\) را ثابت می‌گیریم.
تقسیم دو جریان \eqref{eq:junction-btbt-model} با هندسه و گاف یکسان
می‌دهد
\begin{equation}
 \frac{I_2}{I_1}=
 \frac{1.2}{0.8}\frac{E_{j2}}{E_{j1}}
 \exp\!\left[3\times10^6
 \left(\frac{1}{E_{j1}}-\frac{1}{E_{j2}}\right)\right]
 \approx2.39.
 \label{eq:junction-btbt-ratio-example}
\end{equation}
در این جایگذاری عددی، میدان‌ها برحسب
\(\unit{\volt\per\centi\metre}\) وارد شده‌اند.
برای جریان مطلق، ثابت \(A_b\) و دیگر پارامترهای همان فناوری لازم‌اند.

\paragraph*{\lr{Gate-Induced Drain Leakage}}
\lr{Gate-Induced Drain Leakage (GIDL)} نزدیک ناحیه هم‌پوشانی گیت و
درین رخ می‌دهد. درین بالا و گیت پایین، میدان این ناحیه را تقویت
می‌کنند. در \lr{nMOS}، منفی‌ترکردن گیت در درین بالا می‌تواند این
مسیر را شدیدتر کند. جریان به ناحیه هم‌پوشانی و در نتیجه به عرض
افزاره وابسته است \cite[صفحات ۸۴ و ابتدای ۸۵]{weste}
\cite[اسلاید ۲۸]{nonideal}.
در توصیف نمونه‌های کتاب، سهم این مسیر در
\(|V_{GD}|\le V_{DD}\) معمولاً کوچک است و با راندن گیت بیرون ریل
تغذیه برجسته می‌شود.

برای مثال، در \(V_D=\qty{1}{\volt}\)، تغییر گیت از صفر به
\(-\qty{0.3}{\volt}\)، ولتاژ \(V_{DG}=V_D-V_G\) را از
\(\qty{1}{\volt}\) به \(\qty{1.3}{\volt}\) می‌رساند.
در مدل زیرآستانه، گیت منفی جریان سورس و درین را کاهش می‌دهد، اما
میدان درین و گیت در همین تغییر افزایش می‌یابد.
شکل \secnum{۲.۲۰} کتاب بالا رفتن جریان در سمت گیت منفی را نشان
می‌دهد. برای محاسبه کل نشتی در آن بایاس، همه مسیرهای مربوط به
بایاس و جریان دریافتی از منبع باید در مدل مدار لحاظ شوند.
